\subsection*{Fermi Golden Rule} \index{Fermi Golden Rule}
La formula data dalla regola d'oro di Fermi è tale da permettere di calcolare le probabilità di transizione tra stati quantistici in presenza di una perturbazione dipendente dal tempo. Ovviamente si tratta di un risultato della teoria perturbativa, ed è valido al primo ordine. \\
Si consideri un sistema con una Hamiltoniana $\H_0$ e autostati $\Ket{n}$, a cui si aggiunge una perturbazione dipendente dal tempo $V(t)$. Si vuole valutare l'interazzione tra un livello iniziale $\Ket \vphi$ e una serie di livelli finali densi $\Ket n$, con densità di stati $\rho(E)$. Allora è possibile ricavare la densità di probabilità di transizione in funzione del tempo come:
\[
	\delta \P(t) = \int |\Bra \alpha e^{-\frac{i}{\hslash} \H_0 t} \Ket i|^2 d\alpha = \]\[
	\int \rho(E) |\braket E \psi|^2 dE = \int \rho(E) \frac{1}{\hslash^2} |\Bra E \hat V \Ket {\vphi}|^2F(t, \omega_{i}) dE
\]
dove $F(t, \omega) = \dfrac{4 \sin^2(\omega t/2)}{\omega^2}$ e $\omega_i = (E - E_i)/\hslash$. Per tempi lunghi, $F(t, \omega)$ si comporta come una delta di Dirac, e quindi si può approssimare l'integrale come la sua valutazione al valore $E_i$
\[
	\delta \P(t) = \frac{2 \pi}{\hslash} |V_{E_i, i}|^2 \rho(E_i) t
\]
Che ovviamente è vero solo se l'intervallo di stati di cui si è tenuto conto contiene la $E_i$. La Regola d'Oro di Fermi è data dalla derivata nel tempo della formula appena ricavata:
\[
	\tderiv{}\delta \P(t) = \frac{2 \pi}{\hslash} |V_{E_i, i}|^2 \rho(E_i) = \Gamma_{FGR}
\]

\subsection{Metodo variazionale}
Partiamo da una situazione classica, in cui l'Hamiltoniana è espressa nella sua base come $\H = \sum_n E_n \ketbra n n$, con gli $E_n$ ordinati per energia ($E_n < E_m \Leftrightarrow n < i$). Al ché la media dell'Hamiltoniana uno stato generico è espresso come
\[
	\Bra \psi \H \Ket \psi = \sum_n \braket \psi {E_n} \braket \psi {E_n} < E_0 \sum_n |\braket{\psi}{n}|^2 = E_0 \sfabs{\psi}
\]
Che significa $E_0 < \frac{\Bra \psi \H \Ket \psi}{\sfabs \psi}$.
In sostanza si ha che più si avvicinano le due quantità, più si è vicini alla unica soluzione, che è l'autostato che per energia ha il ground state $E_0$. Vediamo un esempio.

\subsubsection*{Buca infinita di potenziale}
Sia
\[
	V(x) = \begin{cases}
		0 |x| < a\\
		\infty |x| > a
	\end{cases} \]\[
	\psi_0(x) = \frac{1}{\alpha}cos(\frac{\pi x}{2 a})\ ,\ \ \ E_0 = \frac{\hslash^2 \pi^2}{8ma^2}
\]
E variando le $\psi_\lambda(x) = a^\lambda - |x|^\lambda$ risulta
\[
	\avg E = \frac{\Bra {\psi_\lambda} \H \Ket {\psi_\lambda}}{\sfabs {\psi_\lambda}} = \frac{\int_{-a}^a (a^\lambda - |x|^\lambda)(-\frac{\hslash^2}{2m}\sderiv[x]{})(a^\lambda - |x|^\lambda)}{\int_{-a}^a(a^\lambda - |x|^\lambda)^2}
\]

\subsection*{Metodo di Rayleigh-Ritz}
Questo metodo è un'evoluzione del metodo variazionale classico: se sono noti almeno alcuni elementi $\Ket {\chi_i}$ della base ortonormale dell'Hamiltoniana, si può esprimere la $\Ket \psi$ generica come combinazione lineare (complessa) di questi. Al ché si ha:
\[
	\Ket \psi = \sum_{i = 1}^n c_i \Ket {\chi_i} \]\[
	\avg E = \frac{\Bra \psi \H \Ket \psi}{\sfabs \psi} = \sum_{i = 1}^n \sum_{j = 1}^n \frac{c_i^* c_j \Bra {\chi_i} \H \Ket {\chi_j}}{c_i^*c_j}
\]
Si può ora trovare il $c_i^*$ che minimizza $\avg E$:
\[
	\deriv[c_k^*]{E} = 0 = \frac{\sum_j c_j H_{kj}}{\sum_i|c_i|^2} - c_k \frac{\sum_{ij}c_i^*c_jH_{ij}}{(\sum_i|c_i|^2)^2}  
\]
E con questo risultato si può sostituire il peso e ricavare la migliore $\Ket \psi$ come combinazione della base del sottospazio utilizzato.
